
%%%%%%%%%%%%%%%%%%%%%%%%%%%%%%%%%%%%%%%%%%%%%%%%%%%%%%%%%%%%%
\section{模型的评价}

\subsection{模型的优点}
\begin{itemize}[itemindent=2em]
\item \textbf{优点1：几何基础严格。}问题一以半平面求交代替角度区间求交，从表达层面消除了方位角在 $0^\circ$ 与 $360^\circ$ 处的回绕风险；定位区域的有界性、凸性与"直径必在顶点处取到"的性质均有严格论证。直径圆覆盖性以 Thales 逆定理给出充要条件而非仅给出充分条件，结论既完整又可被反例验证，避免了经验性判断。
\item \textbf{优点2：解析与数值互证。}问题二先在垂直布站特例上给出完整闭式解，再在二维决策空间上做数值全局寻优，两条路径在原则上完全独立，所得结论一致，从而排除了解析推导可能出现符号疏漏的风险。同理，问题一的旋转卡壳与暴力枚举、问题三的 Kershner 下界与数值优化也构成互证关系。
\item \textbf{优点3：下界驱动而非经验驱动。}问题三利用 Kershner 覆盖定理确定观测点数下界、利用 Beardwood--Halton--Hammersley 定理确定行程下界，使策略评价有了不可逾越的参照系，把"看起来很快"这类无参照的经验判断替换为"与下界的相对差距"这一可量化指标。
\item \textbf{优点4：判据精确且可判定。}问题四把定向场景下"确保全部清除"这一要求化归为方位角最大间隔小于 $180^\circ$ 的几何判据，给出了零漏测的充要条件，而不是依赖"观测得密一些"这类无法验证的经验做法；由此得到的 $37$ 点冗余网规模是可复核的定量结论。
\item \textbf{优点5：时间口径清晰。}本文明确区分了虚拟世界时间与程序真实运行时间，指出题目真正需要最小化的是前者，后者只构成程序实现效率的约束，避免了以现实时间估算虚拟时间的常见混淆。
\end{itemize}

\subsection{模型的缺点}
\begin{itemize}[itemindent=2em]
\item \textbf{缺点1：观测网构型族受限。}问题三与问题四的观测网均在"中心加同心环"这一构型族内求解。该族具有旋转对称性、便于机器狗沿环巡回，工程上易于实现，但它并非全局最优构型族，所得的 $7$ 点与 $37$ 点结论是族内最优而非全局最优，真实的点数下界可能更低。
\item \textbf{缺点2：分布假设简化。}问题二假定干扰源在可行域内按面积均匀分布，问题三与问题四假定干扰源在目标区域内均匀分布。实际场景中干扰源可能聚集于特定区域，该假设会带来系统性偏差；虽然本文以灵敏度分析讨论了参数变动的后果，但未建立对分布本身的自适应机制。
\item \textbf{缺点3：目标函数单一。}问题二只优化定位精度，未把机器狗的机动成本纳入统一目标，而是以灵敏度分析的形式事后考察。虽然由此得到了"成本权重增大时最优基线向 $54.74^\circ$ 交会角收缩"这一有意义的结论，但未能直接给出精度与时间的联合最优解。
\item \textbf{缺点4：部分参数取经验值。}问题三清除阶段的侧移系数取为 $0.25$，是在"初步定距的位形误差可被后续逼近过程吸收"这一判断下选定的，缺乏与实机参数的标定过程；该系数的取值会对时间模型产生近线性的影响，实际部署前应通过现场试验标定。
\end{itemize}
